\chapter{\ploren{Kod i listingi}{Code and listings}}
\label{app:kod}

\begin{quote}
\itshape
\ploren
  {Załącznik pokazuje sposoby prezentowania kodu źródłowego w pracy. Listing powinien ilustrować omawiany mechanizm, a nie zastępować opisu rozwiązania. Pełny kod projektu należy przechowywać w repozytorium lub dołączyć do materiałów uzupełniających.}
  {This appendix demonstrates methods for presenting source code in a thesis. A listing should illustrate the mechanism being discussed rather than replace the description of the solution. The complete project source should be stored in a repository or included in the supplementary materials.}
\end{quote}

\section{\ploren{Kod wewnątrz zdania}{Code within a sentence}}

Nazwy plików, klas, funkcji i krótkie identyfikatory można zapisywać poleceniem
\verb|\texttt|, na przykład \texttt{moving\_average.py}, \texttt{DataRepository}
i \texttt{sample\_period}. Znaki mające specjalne znaczenie w LaTeX-u trzeba wtedy
poprzedzić ukośnikiem odwrotnym.

Dla krótkiego fragmentu kodu wygodniejsze jest \verb|\lstinline|. Przykładowo funkcję
wywołuje się jako \lstinline[language=Python]|moving_average(data, window=3)|.
Znak po poleceniu \verb|\lstinline| jest ogranicznikiem i nie może wystąpić wewnątrz
prezentowanego fragmentu.

\begin{lstlisting}[style=thesis,caption={Kod LaTeX-a wstawiający identyfikator i fragment kodu w zdaniu},label={lst:kod-w-zdaniu}]
Plik \texttt{moving\_average.py} zawiera funkcję
\lstinline[language=Python]|moving_average(data, window=3)|.
\end{lstlisting}

Listing~\ref{lst:kod-w-zdaniu} przedstawia kod dokumentu, a nie kod programu.
Rozbudowanych instrukcji nie należy umieszczać wewnątrz akapitu, ponieważ utrudniają
czytanie tekstu i mogą przekroczyć szerokość wiersza.

\section{\ploren{Krótki fragment kodu}{Short code excerpt}}

Środowisko \texttt{lstlisting} jest właściwe dla kilku lub kilkunastu wierszy kodu.
Listing~\ref{lst:python-krotki} pokazuje walidację parametru i obliczenie długości
aktualnego okna filtru.

\begin{lstlisting}[style=thesis-python,caption={Walidacja szerokości okna filtru},label={lst:python-krotki}]
def validate_window(index: int, window: int) -> int:
    if window <= 0:
        raise ValueError("window must be positive")
    return min(index + 1, window)
\end{lstlisting}

Opcja \texttt{style=thesis-python} wybiera kolorowanie składni Pythona, numerację
wierszy, łamanie długich linii i podpis umieszczony nad kodem. Listing musi zostać
przywołany przed jego wystąpieniem oraz omówiony po nim. Nie wystarczy stwierdzić,
że „kod przedstawiono poniżej” --- należy wskazać, co jest w nim istotne.

\section{\ploren{Kod w zewnętrznym pliku}{Code stored in an external file}}

Kod wykorzystywany w projekcie powinien znajdować się w rzeczywistym pliku źródłowym.
Polecenie \verb|\lstinputlisting| wczytuje go podczas kompilacji, dzięki czemu listing
nie rozchodzi się z testowaną implementacją. Pełny przykład z Pythona przedstawia
listing~\ref{lst:python-plik}.

\lstinputlisting[
  style=thesis-python,
  caption={Implementacja filtru średniej ruchomej w języku Python},
  label={lst:python-plik}
]{chapters/14-zalacznik-g-kod/code/moving_average.py}

Kod definiuje filtr przyczynowy. Zmienna \lstinline[language=Python]|running_sum|
pozwala uniknąć ponownego sumowania całego okna w każdej iteracji, dlatego złożoność
czasowa wynosi $O(N)$. Instrukcja \lstinline[language=Python]|if __name__ == "__main__"|
tworzy minimalny przykład uruchomienia pliku.

Kod LaTeX-a odpowiedzialny za wczytanie pliku przedstawiono w
listingu~\ref{lst:inputlisting-python}.

\begin{lstlisting}[style=thesis,caption={Wczytanie kodu z zewnętrznego pliku},label={lst:inputlisting-python}]
\lstinputlisting[
  style=thesis-python,
  caption={Implementacja filtru średniej ruchomej w języku Python},
  label={lst:python-plik}
]{chapters/14-zalacznik-g-kod/code/moving_average.py}
\end{lstlisting}

\section{\ploren{Wybór fragmentu pliku}{Selecting a portion of a file}}

W tekście pracy często potrzebny jest tylko fragment implementacji. Opcje
\texttt{firstline} i \texttt{lastline} ograniczają zakres wczytywanych wierszy.
Listing~\ref{lst:python-fragment} przedstawia wyłącznie główną pętlę algorytmu;
numery po lewej odpowiadają numerom w pliku źródłowym dzięki ustawieniu
\texttt{firstnumber=13}, zgodnemu z pierwszym wczytywanym wierszem.

\lstinputlisting[
  style=thesis-python,
  firstline=13,
  lastline=22,
  firstnumber=13,
  caption={Główna pętla filtru średniej ruchomej},
  label={lst:python-fragment}
]{chapters/14-zalacznik-g-kod/code/moving_average.py}

\begin{lstlisting}[style=thesis,caption={Wybór zakresu wierszy z pliku źródłowego},label={lst:inputlisting-zakres}]
\lstinputlisting[
  style=thesis-python,
  firstline=13,
  lastline=22,
  firstnumber=13,
  caption={Główna pętla filtru średniej ruchomej},
  label={lst:python-fragment}
]{chapters/14-zalacznik-g-kod/code/moving_average.py}
\end{lstlisting}

Numery wierszy zmieniają się po edycji pliku, dlatego w opisie merytorycznym lepiej
odwoływać się do nazw funkcji i etapów algorytmu niż utrwalać numery w zwykłym tekście.

\section{\ploren{Przykład w języku MATLAB}{MATLAB example}}

Listing~\ref{lst:matlab-plik} przedstawia tę samą metodę w języku MATLAB. Deklaracja
\texttt{arguments} kontroluje typ i poprawność parametrów, a wcześniejsza alokacja
tablicy \texttt{filtered} ogranicza koszt wielokrotnego powiększania wyniku.

\lstinputlisting[
  style=thesis-matlab,
  caption={Implementacja filtru średniej ruchomej w języku MATLAB},
  label={lst:matlab-plik}
]{chapters/14-zalacznik-g-kod/code/moving_average.m}

\section{\ploren{Przykład w języku C++}{C++ example}}

W listingu~\ref{lst:cpp-fragment} pokazano implementację wykorzystującą typ
\texttt{std::vector<double>}. Parametr wejściowy jest przekazywany jako stała referencja,
a pamięć wyniku jest rezerwowana przed rozpoczęciem pętli.

\lstinputlisting[
  style=thesis-cpp,
  caption={Implementacja filtru średniej ruchomej w języku C++},
  label={lst:cpp-fragment}
]{chapters/14-zalacznik-g-kod/code/moving_average.cpp}

\section{\ploren{Długi kod, wyniki działania i repozytorium}{Long code, program output and repository}}

Listing w pracy powinien być na tyle krótki, aby można go było omówić. Pełnych klas,
plików konfiguracyjnych i wygenerowanego kodu nie należy wklejać wyłącznie w celu
zwiększenia objętości dokumentu. Większy projekt warto udostępnić w repozytorium lub
archiwum wskazanym w informacji o dostępności kodu.

Należy podać wersję języka i najważniejszych bibliotek, sposób instalacji zależności,
polecenie uruchamiające analizę oraz wersję kodu odpowiadającą wynikom w pracy. Dla
Pythona może to być plik \texttt{requirements.txt} lub \texttt{pyproject.toml}, a dla
innych technologii właściwy im plik konfiguracji środowiska budowania.

Wynik działania programu, log lub komunikat terminala nie jest kodem źródłowym. Można
go przedstawić osobnym listingiem bez kolorowania słów kluczowych, wyraźnie opisując,
czy jest to wynik rzeczywistego uruchomienia, czy skrócony przykład.

\begin{lstlisting}[style=thesis,caption={Przykładowy wynik uruchomienia programu},label={lst:wynik-programu},numbers=none,xleftmargin=0.25cm,framexleftmargin=0cm]
$ python moving_average.py
[1.0, 1.2, 1.5, 2.1, 2.6333333333333333]
\end{lstlisting}

\section{\ploren{Najczęstsze błędy}{Common mistakes}}

Do typowych błędów należą: brak odwołania do listingu, podpis pod kodem, wklejanie
nieużywanej implementacji, ręczne kopiowanie kodu zamiast wczytania pliku, brak numerów
wersji i zależności, zbyt długie wiersze, używanie zrzutów ekranu edytora oraz brak
objaśnienia najważniejszych instrukcji. W kodzie przeznaczonym do publikacji nie mogą
znajdować się hasła, tokeny dostępu, prywatne adresy, dane osobowe ani inne informacje
wrażliwe.
